This study is a pilot evaluation rather than a definitive map of model
capabilities. Its primary limitation is reliance on LLM judges. Task-specific
rubrics, multiple judges, leave-family-out masking, and repeated runs improve
auditability but cannot eliminate evaluator bias. Domain-expert validation is
especially important for counseling, tax preparation, tutoring, and
operations research, where professional judgment affects what counts as a
high-quality output.

The augmentation protocol is also intentionally narrow. Each assistant
provides one process-oriented guidance message to a fixed GPT-3.5-Turbo
worker. Real workflows may be multi-turn, allow content-bearing help, or pair
assistants with workers of different capability and model family. Holding the
worker fixed identifies differences among assistants, but assistant rankings
may change with the worker or interaction protocol.

Finally, seven tasks cannot represent the full diversity of professional work.
Although ten independent runs improve stability within this task set, broader
claims require more domains, risk levels, and deliverable formats. Future
validation should test whether the observed cross-role reversals persist with
human workers, domain experts, additional worker models, and authentic
workplace tasks.
